% Source: https://artofproblemsolving.com/community/c3h1285638_dumpling_folding
% Stable identity: aops-topic-1285638-q01
The mass of filling is proportional to the volume of the dumpling, which scales as the cube of the linear size. Doubling the filling multiplies lengths by $\sqrt[3]{2}$.

\begin{asy}
usepackage("amsmath");
size(4cm);
defaultpen(fontsize(10pt));
draw(box((0,0),(1,1)),blue); draw(box((1.6,0),(1.6+1.26,1.26)),red);
label("$5$ cm",(0.5,0),S,blue); label("$5\sqrt[3]{2}$ cm",(1.6+0.63,0),S,red); label("$1$ g",(0.5,0.5),blue); label("$2$ g",(1.6+0.63,0.63),red);
label("\mbox{}",(min(currentpicture,user=true)+max(currentpicture,user=true))/2,black); // keeps the text after the figure black when LaTeX typesets the labels inline
\end{asy}

So the new skins are squares of side $\boxed{5\sqrt[3]{2}\text{ cm}}\approx6.30$ cm. (Their area is multiplied by $\sqrt[3]{4}\approx1.59$, not by $2$.)
